\section{Simulated Fault Scenarios}
\label{sec:possible_anomalies}

The synthetic telemetry generation framework enables the reproduction of several fault conditions commonly reported in industrial AGV and AMR systems. Unlike isolated sensor malfunctions, these anomalies affect multiple telemetry channels simultaneously due to the physical coupling between electrical, mechanical, and thermal subsystems. Consequently, they produce characteristic multivariate signatures that can be used to validate anomaly-detection algorithms.

The fault scenarios implemented in the simulator include mechanical stall conditions, wheel slippage events, gearbox and bearing degradation, stator winding faults, battery ageing effects, and inverter phase failures. These scenarios were selected because they represent common failure modes in electromechanical drive systems and generate characteristic telemetry patterns suitable for validating predictive maintenance and anomaly detection techniques.

Based on the electromechanical relationships presented in Section~\ref{sec:synthetic_dataset_generation}, each failure mode can be associated with a specific combination of telemetry deviations. Table~\ref{tab:diagnostic_matrix} summarizes the simulated fault conditions together with their expected telemetry signatures, underlying physical mechanisms, and diagnostic conditions.

\begin{table}[htbp]
\centering
\caption{Multivariate Telemetry Diagnostic Rules for AGV/AMR Powertrains}
\label{tab:diagnostic_matrix}
\small
\begin{tabular}{p{2.8cm}p{3.5cm}p{4.2cm}p{4.0cm}}
\toprule
\textbf{Fault State} & \textbf{Telemetry Pattern} & \textbf{Physical Mechanism} & \textbf{Diagnostic Condition Rule} \\
\midrule
\textbf{Mechanical Stall / Rotor Lock} & $I_{\mathrm{m}} \rightarrow I_{\mathrm{stall}}$, $v = 0$, $V_{\mathrm{m}} \rightarrow V_{\mathrm{bat}}$, $\frac{\mathrm{d}T_{\mathrm{m}}}{\mathrm{d}t} \gg 0$, $V_{\mathrm{bat}} \downarrow\downarrow$ & Mechanical blockage prevents rotor rotation ($\omega_{\mathrm{m}}=0$), causing the back-EMF term to collapse ($E_b = 0$) and producing a large current surge. & $(I_{\mathrm{m}} > 2.5 \cdot I_{\mathrm{rated}})\land(v < 0.01)\land(v_{\mathrm{ref}} > 0.1)$ \\
\addlinespace
\textbf{Wheel Slippage} & $v_{\mathrm{enc}} \gg v_{\mathrm{ground}}$, $I_{\mathrm{m}} \downarrow$, $V_{\mathrm{m}} =$ Nominal, $T_{\mathrm{m}} =$ Low & Loss of traction reduces load torque ($\tau_L \rightarrow 0$) while the motor continues rotating. & $\left( \frac{v_{\mathrm{enc}}-v_{\mathrm{ground}}} {v_{\mathrm{ground}}} \right)>0.15 \land (I_{\mathrm{m}} < I_{\mathrm{base}}) $ \\
\addlinespace
\textbf{Gearbox / Bearing Wear} & $I_{\mathrm{m}}\uparrow$, $v=v_{\mathrm{ref}}$, $T_{\mathrm{m,ss}}\uparrow$, $V_{\mathrm{bat}} =$ Nominal & Additional friction generated by bearing degradation or gearbox wear increases mechanical losses. & $(\Delta I_{\mathrm{m,mean}} > 20\%) \land (v=v_{\mathrm{ref}})$ \\
\addlinespace
\textbf{Stator Inter-Turn Short} & $I_{\mathrm{m,ripple}}\uparrow$, $T_{\mathrm{m}}\uparrow\uparrow$, $V_{\mathrm{m}} =$ Unbalanced, $v \approx v_{\mathrm{ref}}$ & Partial insulation breakdown produces circulating currents and localized heating. & $\mathrm{THD}(I_{\mathrm{m}}) > \mathrm{Threshold} \land \frac{\mathrm{d}T_{\mathrm{m}}}{\mathrm{d}t} > \mathrm{Baseline} $ \\
\addlinespace
\textbf{Battery SOH Degradation} & $I_{\mathrm{m}} =$ Nominal, $v =$ Nominal, $V_{\mathrm{bat}}\downarrow\downarrow$ during acceleration & Internal battery resistance increases due to electrochemical aging, producing excessive voltage sag. & $\frac{\Delta V_{\mathrm{bat}}} {\Delta I_{\mathrm{bat}}} = R_{\mathrm{int}} > 1.5 \cdot R_{\mathrm{int,nom}} $ \\
\addlinespace
\textbf{Inverter Phase Loss} & $I_{\phi A}=0$, $I_{\phi B,C}\uparrow\uparrow$, $v\ll v_{\mathrm{ref}}$, $T_{\mathrm{inv}}\uparrow\uparrow$ & Open-switch failure or phase disconnection disrupts torque generation and current balance. & $\min(I_\phi)=0 \land \max(I_\phi)>I_{\mathrm{rated}} $ \\
\bottomrule
\end{tabular}
\end{table}
